\subsection*{Stage 004.006: Near-throat projection scale}
\label{subsec:stage004006-projected-maxwell-near-throat}

\paragraph{Source and audit.}
This substage corresponds to
\StageFile{scripts/moving_throat_pde_stage004_006_projected_maxwell_near_throat_sympy_audit.py}.
It distinguishes symmetric interior averaging from one-sided mouth averaging.

\paragraph{Interior versus mouth.}
For a symmetric interior kernel, the first finite-width correction to a smooth
profile is even and begins at \(O(\sigma^2)\).  At a mouth with one-sided
coordinate \(s\ge0\),
\begin{equation}
\label{eq:stage004006-mouth-kernel}
W_\ell(s)=\frac1\ell w(s/\ell),
\qquad
\int_0^\infty w(u)\,du=1,
\end{equation}
the expansion is
\begin{equation}
\label{eq:stage004006-mouth-expansion}
\langle X\rangle_\ell
=X(0)+\ell\mu_1X'(0)+O(\ell^2),
\qquad
\mu_1=\int_0^\infty u\,w(u)\,du .
\end{equation}

\paragraph{Output.}
Stage~004.006 explains why mouth-local projected EM data enter at first order
in the mouth scale, while symmetric interior thickness effects enter later.
